\documentclass[11pt,a4paper]{article}
\usepackage{gtu-el-yazisi}
\begin{document}
\notbilgi{Alan Teorisi}{Ders 1 -- Kaynaklar ve temel varsayımlar}{Ders1.pdf}

\sayfa{1}
\section{1. ders}
\underline{1. kitap}

\underline{2. kitap :} Elektromanyetik Alanteorisinin çözümlü problemleri, 2019

\hspace{1.5em}Prof.dr G.Uzgören, A.Büyükaksoy, A.Alkumru\hfill , 2023

\underline{Ek kaynak:} 1: David K CHENG çeviri: Mühendislik Elektromanyetiğin temelleri

\hspace{1.5em}Prof.dr.Adnan Kösal, Birsen Saka, 2. Baskı 2020

\hspace{1.5em}2. David K CHENG \quad Elelectromagnetics Fields \& Waves

\medskip
\noindent
\begin{tabular}{@{}l@{\quad}l@{\quad}l@{\quad}l@{}}
  2 ödev ve bir bilgisayar projesi & & & (Her biri 5'er puan) \\
  1 tane ara sınav & 23 kasım 2023 & (Perşembe ders saati) & (Vize 35 Puan) \\
  Final & 25 ocak 2024 & (\hspace{2em}''\hspace{2em}) & (Final sınavı 50 puan) \\
\end{tabular}

\section{1. Ders Notu}
\noindent
\begin{tabular}{@{}l@{\qquad}l@{\qquad}l@{}}
  İki etkileşim türü vardır : & 1. Mekanik & 2. Elektriksel \\
  $(\vec{r},t)$, \ $\vec{r}\to x,y,z$ & zaman, hız, ivme, konum & \\
\end{tabular}

\medskip
\noindent
\begin{minipage}[t]{0.47\linewidth}
  \textbf{\underline{1. Temel Varsayımlar}}
\end{minipage}\hfill
\begin{minipage}[t]{0.47\linewidth}
  \textbf{\underline{İlk kavramlar ve Sonuçlar}}
\end{minipage}

\sekilnotu[1]{Bu sayfa basılı materyal içeriyor (ör. ders slaytı); dijitalleştirilmedi. (Yapıştırılmış basılı tablo: ``Temel varsayımlar'' (11 madde) ve ``İlk kavramlar ve sonuçlar'' sütunları. El yazısı işaretler: 1.--4. maddeler sarı vurgulu ve solda mavi köşeli parantezle ``Mekaniksel'' diye işaretlenmiş; 11. maddenin numarası sarı vurgulu; sağ sütunda ``Yüzeysel yük ve yüzeysel akım, sınır koşulları'' satırının yanına mavi el yazısıyla ``, dipol dağılımı, ayırt koşulları'' eklenmiş.)}

% Sayfa 2: bos (sadece cizgili kagit), aktarilmadi.
\end{document}
